% 第 3 题频率分布直方图：按固定源 PDF 矢量路径实测重绘（1 pt → 0.03515 cm）。
% 原点取轴交点（源页 (132.48, 466.41) pt），y 向上。组距 2.5 小时 = 21.26 pt；
% 横轴原点右侧带折断符号（锯齿 x 4.25–17.00、振幅 ±2.13，照源复刻）；
% 纵轴刻度自下而上 0.02/0.04/0.08/0.10/0.16（无 0 标签）；柱高 14.17/70.84/113.34/56.67/28.33；
% 5 条水平虚线自纵轴右行、止于对应柱左边缘；纵轴标签「频率／组距」为竖排分数式。
\begin{tikzpicture}[>=stealth,every node/.style={font=\small},x=0.03515cm,y=0.03515cm]
% 坐标轴（纵轴至箭头尖端 134.6；横轴带原点折断符号）
\draw[->] (0,0)--(0,134.6);
\draw (0,0)--(4.25,0);
\draw (4.25,0)--(5.32,2.13)--(7.44,-2.13)--(9.57,2.13)--(11.69,-2.13)--(13.82,2.13)--(15.94,-2.13)--(17.00,0);
\draw[->] (17.00,0)--(170.0,0);
% 柱（高 0.02/0.10/0.16/0.08/0.04，组距 21.26）
\foreach \i/\h in {0/14.17, 1/70.84, 2/113.34, 3/56.67, 4/28.33} {
  \draw ({21.25+21.26*\i},0) rectangle ({21.25+21.26*(\i+1)},{\h});
}
% 水平虚线引导线（自纵轴至对应柱左缘）
\draw[dash pattern=on 2.99pt off 2.99pt] (0,14.17)--(21.25,14.17);
\draw[dash pattern=on 2.99pt off 2.99pt] (0,70.84)--(42.51,70.84);
\draw[dash pattern=on 2.99pt off 2.99pt] (0,113.34)--(63.76,113.34);
\draw[dash pattern=on 2.99pt off 2.99pt] (0,56.67)--(85.01,56.67);
\draw[dash pattern=on 2.99pt off 2.99pt] (0,28.33)--(106.26,28.33);
% 纵轴刻度标签（虚线高度处，纵轴左侧）
\node[left] at (-2.4,14.17) {\(0.02\)};
\node[left] at (-2.4,28.33) {\(0.04\)};
\node[left] at (-2.4,56.67) {\(0.08\)};
\node[left] at (-2.4,70.84) {\(0.10\)};
\node[left] at (-2.4,113.34) {\(0.16\)};
% 横轴刻度标签（组边界，居中于刻度）
\node[below] at (0,-2.6) {\(0\)};
\foreach \t/\x in {17.5/21.25, 20/42.51, 22.5/63.77, 25/85.03, 27.5/106.29, 30/127.55} {
  \node[below] at (\x,-2.6) {\(\t\)};
}
% 纵轴标签「频率／组距」竖排分数式（照源：频率在上、横线、组距在下）
\node[anchor=west,inner sep=0] at (4.90,142.15) {频率};
\draw (4.89,133.91)--(25.89,133.91);
\node[anchor=west,inner sep=0] at (4.90,127.84) {组距};
% 横轴标签「自习时间/小时」（横轴上方、右端；基线在横轴上方 6.32）
\node[anchor=base west,inner sep=0] at (136.0,6.32) {自习时间/小时};
\end{tikzpicture}
